\clearpage
\section*{Review, notation, and sources}\label{sec:review}
\subsection*{Explain the complete modeling path}
Given a new text sample, you should be able to identify what preprocessing
preserves, how tokens are formed, which context each prediction uses, how
its probability is normalized, and how a learning signal reaches parameters.
The following questions connect the sections.

\begin{tabularx}{\textwidth}{@{}l Y@{}}\toprule
Topic & A question to answer without running code\\\midrule
Tokenization & Why can byte coverage coexist with poor compression for a language?\\
Counts & How do unseen histories differ from unseen continuations of a known history?\\
Evaluation & What changes in perplexity when the token unit changes?\\
Embeddings & What distinguishes pair discrimination from normalized next-token likelihood?\\
Neural LMs & Why cannot a wider token table alone distinguish two prefixes ending in the same token?\\
Training & What does successful tiny-batch fitting establish, and what requires held-out evidence?\\
Recurrence & How can the forward context exceed the truncated backward horizon?\\
LSTM & Which derivative does a product of forget gates describe?\\\bottomrule
\end{tabularx}

\subsection*{Notation to keep consistent}
$x_t$ is a token ID; a state after consuming $x_t$ predicts $x_{t+1}$.
$K_{\mathrm{in}}$ counts input rows and $K$ counts possible outputs.
$B,T,d,m,H$ denote batch size, input positions, embedding width, fixed
context length, and recurrent/hidden width. $E$ is the input table;
readout matrices are identified beside their equations. In the LSTM cell,
the subscript $o$ names its output gate rather than the vocabulary readout.
$\log$ means natural logarithm; $\log_2$ measures bits. MiB means $2^{20}$ bytes.

\subsection*{Course sources and numerical provenance}
The source baseline is repository commit \texttt{1adbaa5}. The
\href{https://github.com/baojian/llm-26-fall/tree/1adbaa5442a567cbd2ab65035b8e485778dd46a9/slides/lecture-01}{L01},
\href{https://github.com/baojian/llm-26-fall/tree/1adbaa5442a567cbd2ab65035b8e485778dd46a9/slides/lecture-02}{L02},
\href{https://github.com/baojian/llm-26-fall/tree/1adbaa5442a567cbd2ab65035b8e485778dd46a9/slides/lecture-03}{L03}, and
\href{https://github.com/baojian/llm-26-fall/tree/1adbaa5442a567cbd2ab65035b8e485778dd46a9/slides/lecture-04}{L04}
directories contain the slides, teaching plans, and classroom notebooks.
The revised notes expand recurrent models and LSTMs at the instructor's
request and place attention in Lecture 05; the older L04 deck has a different boundary.

\begin{itemize}
\item In the lecture directories, F07 redraws L01's
\path{assets/heldout-chinese.json}; F15 uses L03's \path{assets/counts-ppmi.json}.
\item F23 redraws the two distinct \texttt{tiny-lm-loss.json} assets in
L03 and L04. Their captions state models, batches, seeds, and update budgets.
\item The L02 loop numbers come from \texttt{self-training-loop.json}.
Full settings and data fingerprints are in \texttt{lecture02-results.json}
and its asset README; upstream dataset revisions were not recorded for the original download.
\item Other numerical illustrations are exact arithmetic or explicitly
chosen toy values. No new model training is reported by these notes.
\end{itemize}

The next four pages give the selected 15-paper reading route. Supporting
sources for gradient clipping and the later forget gate are linked in
Sections 7--8. The subsequent Transformer paper is
\href{https://arxiv.org/abs/1706.03762}{Vaswani et al.\ (2017)}; its
architecture belongs to the next part of the course.
